% ------------------------------------------------------------------------------
% Slide 9: Δακτύλιος Chord & Learned Index (RMI)
% ------------------------------------------------------------------------------
\begin{frame}{Δακτύλιος Chord \& Learned Index (RMI)}
  \begin{columns}[c]
    \begin{column}{0.50\textwidth}
      \textbf{Κατανεμημένος Δακτύλιος LEAD:}
      \begin{itemize}
        \item 100 Virtual Nodes (vnodes) ανά κόμβο.
        \item Δρομολόγηση $O(\log N)$ hops με finger tables.
      \end{itemize}

      \vspace{0.1cm}
      \textbf{Ευρετήριο RMI (Recursive Model Index):}
      \begin{itemize}
        \item Αντικατάσταση B-Tree με μοντέλο 2 επιπέδων (Root Linear $\to$ 256 Leaf Linear Models).
        \item \textbf{Αποτέλεσμα:} Χρόνος ανάκτησης $\approx 41\text{ ns}$ ($1.81\times$ ταχύτερο του B-Tree).
        \item \textbf{$3\times$ χαμηλότερη μνήμη RAM} ($785\text{ KB}$ vs $2.34\text{ MB}$).
      \end{itemize}
    \end{column}

    \begin{column}{0.48\textwidth}
      \centering
      \includegraphics[width=\textwidth,height=0.65\textheight,keepaspectratio]{fig3_rmi_vs_btree_benchmark.png}
      \par\vspace{0.05cm}{\centering\scriptsize RMI vs std::BTreeMap Latency.\par}
    \end{column}
  \end{columns}
\end{frame}
